\section{模型假设与符号说明：让光学条件可追溯}

本章把每项简化固定在它实际发生的建模位置，并给出可复核的检验。模型将空气、外延层和衬底分成三层，外部入射角取 $10^\circ$ 与 $15^\circ$；三层结构用于界面反射、透射和纵向传播因子的表达，分别由场展开互验、被动传播支检查和双角重估核对。
% src:交接/数据档案.json:总览.每块晶圆入射角度
% src:交接/假设台账_问题1.json:0.假设;0.检验结果
% 依据:交接/建模笔记_问题1.md:2. 几何、波数和传播因子

三问依次使用首个返回场、双角共享相位与完整往返场；透明近似、偏振权重和有效损耗的适用位置由假设表逐项限定。
% src:交接/假设台账_问题1.json:1.假设;1.检验结果
% src:交接/假设台账_问题2.json:3.假设;3.检验结果
% src:交接/假设台账_问题3.json:3.假设;3.检验结果;4.假设;4.检验结果;5.假设;5.检验结果

单位约定贯穿三问：反射率先按比例进入计算，波数以 $\mathrm{cm}^{-1}$ 记录，厚度在相位中以厘米计算，输出时再换算为微米。厚度沿 $8\,\mu\mathrm{m}\rightarrow0.0008\,\mathrm{cm}\rightarrow8\,\mu\mathrm{m}$ 回算一致，说明波数与厚度相乘的相位无量纲；下面两小节将把这些生效位置展开为假设表，并用统一符号承接三问的场模型。
% src:求解/问题1/结果/模型验证.json:单位换算.输入厚度_微米;单位换算.厚度_厘米;单位换算.回换厚度_微米
